% TikZ 图：磁场叠加
% Day 21：磁场与安培力
\begin{tikzpicture}[scale=0.9,>=Latex]
  \draw[thick] (-2.1,0) circle (0.20);
  \fill[red] (-2.1,0) circle (0.055);
  \node[font=\small] at (-3.05,-0.05) {$I_1$ 向外};
  \draw[blue,thick,->] (-2.1,0.65) arc (90:380:0.65);

  \draw[thick] (2.1,0) circle (0.20);
  \draw[red,thick] (1.98,-0.12) -- (2.22,0.12);
  \draw[red,thick] (1.98,0.12) -- (2.22,-0.12);
  \node[font=\small] at (3.05,-0.05) {$I_2$ 向里};
  \draw[blue,thick,->] (2.1,0.65) arc (90:-200:0.65);

  \fill[red] (0,1.15) circle (2pt) node[above=2pt] {$P$};
  \draw[->,blue,thick] (0,1.15) -- (-0.82,1.72) node[left] {$B_1$};
  \draw[->,blue,thick] (0,1.15) -- (0.82,1.72) node[right] {$B_2$};
  \draw[->,red,very thick] (0,1.15) -- (0,2.18) node[above] {$B_{\text{合}}$};

  \node[font=\small,align=center] at (0,-1.10)
    {$|I_1|=|I_2|$，且 $P$ 到两导线距离相等};
  \node[font=\small,align=center] at (0,-1.52)
    {先分别判断 $\boldsymbol B_1,\boldsymbol B_2$，再按平行四边形法则叠加};
\end{tikzpicture}
